\section {Существующие LSM-деревья}

Данный раздел посвящён тому, какие вариации внутреннего устройства имеет LSM-дерево, для чего нужны те или иные модификации, а так же в каких реальных системах управления базами данных они используются. 

Сначала давайте обсудим некоторые допущения, которые будем использовать в рассуждениях ниже.
Источник \cite{luo2020lsm} утверждает, что для LSM-дерева под стабильной нагрузкой (когда количество слоёв не меняется), перфоманс записи оптимален, когда для любого $L$ максимальный количество ключей на $L + 1$ уровне в $T$ раз больше, чем максимальное количество ключей на $L$ уровне. 

% \subsection{Дисковые компоненты}

% В LSM-деревьях может быть по разному устроено хранение дисковых компонент
% \begin{paracol}{2}
%   %TODO сделать не так категорично
%   \begin{center}\bfseries $B^+$ -дерево\end{center}
%   Используя данный метод мы во многом теряем два главных плюса LSM - дерева -- мало чтений со случайного места в диске и отказоустойчивость, потому что LSM - дерево с SSTable внедрялось, как альтернатива $B^+$-дерева, с такими свойствами. Так же амортизированно получается более долгая вставка и удаление. Единственный очевидный плюс: мы можем не хранить разные версии одного ключа, проверяя его наличие в компонентах LSM-дерева, и явно перезаписывать/удалять. данная идея изначально кажется странной, но авторы вот этой статьи \cite{luo2020lsm} вскользь упоминают данную реализацию, поэтому она стоит упоминания.
% \switchcolumn
%   \begin{center}\bfseries SSTable\end{center}
%   В такой реализации мы получаем от 
%   LSM дерева все те качества, которые и хотели получить: мало случайных I/O запросов при записи новых элементов, и слиянии существующих компонентов 
%   LSM - дерева. Отказоустойчивость: корректность записи одного элемента на диск довольно просто отследить и исправить, по сравнению с перестройкой $B^+$-дерева. С другой стороны, любые слияние таблиц происходят в дополнительной памяти, поэтому данные в таблицах никогда не будут испорчены.
% \end{paracol}
%   Так же в этой статье упоминаются различия в 

\subsection{Политика слияния}

LSM-деревья различаются между собой по политике слияния дисковых компонент. \cite{luo2020lsm} На практике обычно используют следующие две политики, каждая заточена под свой преобладающий тип запросов:

\begin{paracol}{2}
  \begin{center}\bfseries Leveling Merge Policy\end{center}

  Данная политика хранит компоненты следующим образом:

  На каждом уровне будет хранится ровно одна компонента.

  Максимальный объем компоненты на уровне $L$ хранит в $T$ раз больше ключей, чем компонента на уровне $L - 1$. Тогда компонента на слое $L$ сольется с компонентой на $L - 1$  слое много раз (в точности $T$) перед тем, как переполнится и начнёт сливаться с компонентой на слое $L + 1$.

  Как итог, Leveling Merge Policy хранит "мало" компонент (по сравнению со второй политикой), благодаря чему, показывает лучший перфоманс на запросах чтения. 

  
\switchcolumn
  \begin{center}\bfseries Tiering Merge Policy\end{center}

  Пусть вместимость $L-1$ слоя $P$ - ключей, тогда
  эта политика, на каждом слое $L - 1$ будет хранится не более $T$ компонент, каждая размера $\frac{P}{T}$ с возможно пересекающимся диапазоном ключей. Когда на слое наберётся $T$ заполненных компонент, сольем их в одну компоненту и положим на следующий уровень. Действительно, такое поведение будет корректным, потому что компонента на $L$ слое вместит в себя $T$ раз больше ключей, т.е. $P$, и максимальная вместительность уровня $L$ в точности $P \cdot T$, что в $T$ раз больше, чем на предыдущем уровне. И заполненный $L-1$ уровень станет ровно одной компонентой на $L$ слое.

  Эта политика, напротив, оптимизирована для запросов записи, за счёт снижения числа тяжёлых слияний, когда структуре приходится сливать старшие уровни.


\end{paracol}




Вообще говоря, классической политикой для LSM-дерева является именно Leveling Merge Policy, но параллельно LSM-дереву,  Jagadish и др. в своей статье \cite{jagadish1997incremental} представили очень похожую на LSMT структуру, оптимизированую под запросы записи, в которой применили "stepped-merge policy". Которая в последствии превратилась в Tiering Merge Policy, используемую в современных СУБД.


\subsection{Some Well-Known Optimizations}

\epigraph{Название данной секции дословно ссылается на название секции из статьи: "LSM-based Storage Techniques: A Survey" \cite{luo2020lsm} }

\subsubsection{Фильтр Блюма и подобные}

Фильтр Блюма - вероятностный алгоритм, для проверки принадлежности элемента множеству, гарантирующий отсутствие ложно-отрицательных срабатываний.

Устройство фильтра: есть фиксированное $m$ - количество битов и $k$ и $k$ хэш функций $h_1, \ldots, h_k $, из некоторого семейства, $h_i:\mathcal{K}\rightarrow \{0, \ldots, m-1 \}$. Положим массив $A$ битов размера $m$. Когда в множество добавляется новый ключ $K$, положим $\forall i \in \{1, \ldots, k\} A[h_i(K)] := 1$. Проверка принадлежности ключа множеству проводится следующим образом: необходимо проверить, что $\forall i A[h_i(K)] == 1$.


Фильтр Блюма позволяет проверить отсутствие ключа в SSTable, перед началом дорого (с точки зрения числа I/O запросов) поиска ключа в самой таблице. В свою очередь, ложно-положительные срабатывания не влияют на корректность ответа, только лишь добавляют расходы на поиск внутри самой компоненты. Вероятность такого  ложно-положительного срабатывания фильтра не превосходит $(1 - e ^ {-kn / m})^k$, тогда оптимальное количество хэш  функций для минимизации ошибки достигается при $k:=\frac{m}{2}\ln 2$, что на практике оказывается ~10 и позволяет достичь $1\%$ ложно-положительных срабатываний. А так как фильтр Блюма очень маленький и может быть загружен в оперативную память, накладные расходы на $I/O$ запросы к диску во многих сценариях получается существенно сократить.

Собственно, идея о том, что с помощью некоторой небольшой вероятностной структуры можно с хорошей точностью быстро проверить отсутствие ключа (если ключа нет в таблице, то структура должна гарантированно это сообщить) в одной компоненте, находит применение во многих современных СУБД. LevelDB использует встроенный фильтр Блюма - NewBloomFilterPolicy \cite{leveldbgit}. А RocksDB позволяет выбирать фильтр: NewBloomFilterPolicy, NewRibbonFilterPolicy (Ribbon Filter - некоторый аналог фильтра блюма, который занимает меньше памяти).
 \cite{rocksdbgit}.


\subsubsection{Partitioning}

%TODO я не хочу это писать


\subsection{LevelDB}

\subsection{RocksDB}

\subsection{Apache Cassandra / ScyllaDB}

\subsection{Apache HBase}

\subsection{CockroachDB}

\subsection{WiredTiger}

